\section{What Goes In Each Crate, and Why}
\label{sec:what-goes-in-each-crate}

The PAC/HAL/BSP pattern from Section~\ref{sec:pac-hal-bsp-architecture} says the layers exist.
It does not say where one ends and the next begins, and every project fills that gap
differently. That gap is what you face the first time you add support for a new MCU or a new
board — which is the situation this book is written for, and the one it should leave you
equipped to handle.

This section argues the choice, compares it with what four mature Rust embedded ecosystems do,
gives a decision procedure you can apply to hardware that is not in this book, and states the
rules that keep the answer from drifting.

\subsection{The Question, Stated Sharply}

Three pieces of code make the ambiguity visible. All three are real, from this project:

\begin{enumerate}
    \item \textbf{A pin pair.} The IMU on this board answers on \texttt{SDA = PB2},
          \texttt{SCL = PB3} — the opposite of what the vendor's pin map declares. Whose
          knowledge is that?
    \item \textbf{A power gate.} The same sensor does not answer until \texttt{PD5} is driven
          high. Is that a driver detail or a board detail?
    \item \textbf{An algorithm.} A fall classifier takes acceleration samples and decides
          whether the device fell. It has no registers, no pins and no clock. Where does it
          live so that it stays testable?
\end{enumerate}

A HAL that answers (1) and (2) knows too much: every new board becomes a fork of the HAL. An
example that answers them re-derives them, and the same investigation gets repeated by the
next person — the failure this project has already paid for once.

\subsection{Three Candidate Layouts}

\paragraph{Layout A — no BSP at all; wiring lives in the examples.}
This is the smallest possible structure: the HAL offers drivers, and each example says which
pins to use. It is what the STM32 and ESP32 ecosystems do, and it is defensible when the point
of the example is to teach the peripheral.

Its cost appears the moment the wiring is not obvious. This board's IMU needed a swapped
\texttt{SDA}/\texttt{SCL} and a power gate before it answered anything; putting that in every
example means every example can get it wrong independently, and the knowledge is only as good
as the most recently written one.

\paragraph{Layout B — one BSP crate, features for the optional hardware.}
A single \codeinline{xiao-}\allowbreak\codeinline{mg24-bsp} with features
\codeinline{sense} and \codeinline{expansion-base}. It looks like the cheapest correct answer, and it has a specific
trap that the nRF ecosystem documents in its own repository: \emph{cargo unifies features
across a workspace}, so a build that enables the base for one target can silently enable it
for another. More importantly, the API lies: a \codeinline{ExpansionBase} type that exists in a
build without the base has to answer ``no OLED'' at runtime, and every method that could have
been a compile error becomes an \codeinline{Option}.

\paragraph{Layout C — one BSP crate per board, and a feature for a variant of one board.}
The module gets its own crate (LED on both variants; IMU and microphone behind
\codeinline{sense}), the Expansion Base gets its own crate that depends on the module's, and
the dependency arrow points one way. A user who does not own the base never compiles it, and
never sees an \codeinline{Option} that only exists because two boards were welded together.

\subsection{What Four Ecosystems Actually Do}

Read from the repositories, not from memory:

\begin{table}[H]
\centering
\footnotesize
\begin{tabular}{@{}>{\raggedright\arraybackslash}p{2.3cm}>{\raggedright\arraybackslash}p{2.6cm}>{\raggedright\arraybackslash}p{2.4cm}>{\raggedright\arraybackslash}p{3.4cm}>{\raggedright\arraybackslash}p{3.0cm}@{}}
\hline
\textbf{ecosystem} & \textbf{chip selection} & \textbf{BSP crates} & \textbf{where examples live} & \textbf{where wiring is written} \\
\hline
\textbf{rp-hal} (RP2040/RP2350) & one HAL per chip & \textbf{yes}, one per board,
  in a \emph{separate repository} (\codeinline{rp-hal-boards/boards/}) & in each BSP, plus
  separate example crates & in the BSP, once \\
\textbf{nrf-hal} & one HAL per chip (\codeinline{nrf52840-hal}, \codeinline{nrf9160-hal}, \dots)
  + \codeinline{nrf-hal-common} & no & in the HAL repo, board chosen by feature & in the
  examples \\
\textbf{stm32f4xx-hal} & one HAL per family & no & in the HAL crate &
  in the examples (pin numbers inline) \\
\textbf{esp-hal} & one crate, chip by feature & no & one \codeinline{examples/} grouped by
  capability & in the examples \\
\hline
\end{tabular}
\caption{How four Rust embedded ecosystems split PAC, HAL, BSP and examples. Verified from the
repositories in November 2026.}
\label{tab:ecosystem-layouts}
\end{table}

Two details in that table are worth more than the table.

First, \textbf{rp-hal puts its BSPs in a different repository from the HAL}. That is stronger
than a folder: the HAL's continuous integration cannot depend on a board, because the board is
not present. The BSP crates describe themselves as a HAL for one board that re-exports the
chip HAL — a facade, exactly as this book calls it.

Second, \textbf{nrf-hal ships a workspace template instead of a workspace}, and says why:
\emph{``Since Cargo resolves all features for all crates in the workspace together, it is not
possible to have a workspace with all the crates together, since they enable incompatible
features.''} That sentence is the argument against Layout B in the ecosystem's own words.

\subsection{The Decision, and the Rule Behind Each Part}

\begin{table}[H]
\centering
\footnotesize
\begin{tabular}{@{}>{\raggedright\arraybackslash}p{3.4cm}>{\raggedright\arraybackslash}p{5.2cm}>{\raggedright\arraybackslash}p{5.2cm}@{}}
\hline
\textbf{crate / folder} & \textbf{what goes in} & \textbf{what must never go in} \\
\hline
\codeinline{efr32mg24-pac} & registers, generated from the SVD & anything hand-written that
  a generator could produce \\
\codeinline{efr32mg24-hal} & drivers: which bytes, in which order, with which timing & pins,
  addresses, power gates, board names \\
\codeinline{xiao-mg24-bsp} & the module: LED; behind \codeinline{sense}, the IMU bus, its
  power gate and the microphone & any device that is not on the module \\
\codeinline{xiao-expansion-base-bsp} & the base: OLED, BME68x, RTC, buzzer, button, Grove,
  microSD, and the base's own bus & anything that exists without the base \\
\codeinline{libraries/} & domain policy with no hardware at all & any dependency on the
  hardware stack, direct or transitive \\
\hline
\end{tabular}
\caption{Ownership. The right-hand column is what the gate checks; the left two are what a
review checks.}
\label{tab:crate-ownership}
\end{table}

Three rules follow, and each one exists because of a specific failure:

\begin{keytakeaways}
    \item \textbf{A BSP describes one piece of plastic.} Two boards that stack are two
          crates, because a user may own only one of them.
    \item \textbf{A variant of one board is a feature, not a crate.} The Sense module is the
          same PCB with more parts fitted; a separate crate would duplicate the board init
          and the LED to add two sensors.
    \item \textbf{A dependency arrow never points at a board.} The base depends on the
          module; the module does not know the base exists. This is the property that makes
          ``does it work without the base?'' answerable by the compiler instead of by a test
          matrix.
\end{keytakeaways}

\subsection{Where the Algorithm Goes}

The classifier deserves its own sentence because it is where layering usually breaks. It has
no hardware, but the temptation is to put it ``near'' the driver that produces its input — in
the same crate as the IMU, or in the example that uses it.

Both are wrong for the same reason: the moment it shares a crate with anything that needs a
target triple, its tests need a target triple, and the tests are what keep it honest. In this
project it is a crate with \textbf{zero dependencies}, and the gate fails if its dependency
tree ever contains a hardware crate:

\begin{minted}{bash}
$ cargo tree -p motion-detection | grep efr32mg24
$ echo $?
1        # nothing from the hardware stack: the rule holds
\end{minted}

The contract is worth quoting, because it is the whole testability argument in one line:

\begin{minted}{rust}
// [i32; 3] of milligravity, and the interval between samples in microseconds.
pub fn update_us(&mut self, accel_mg: [i32; 3], dt_us: u32) -> Option<Event>
\end{minted}

No sensor type, no bus, no clock. Which is why nine tests run in under a second on a laptop,
and why the same code can be replayed against a recording at ten times real speed.

\subsection{And Where the Examples Go}

The layout of the examples follows the same principle as the crates, with one addition for
hardware that stacks: \textbf{an example lives in the crate of the largest hardware
requirement it has, and states that requirement in its first comment.}

\begin{minted}{text}
crates/xiao-mg24-bsp/examples/            runs with the module alone on the desk
crates/xiao-expansion-base-bsp/examples/  needs the module (Sense) and the base
crates/efr32mg24-pac/examples/            chip level: any pin is legitimate
\end{minted}

The alternative — every example in one folder — looks tidier and is worse: the crate of the
module would have no examples, and an API that nothing builds is an API that nothing checks.
A crate with hardware features and zero examples should fail a gate, not pass it quietly.

\subsection{A Decision Procedure for Hardware Not In This Book}

The four questions below decide the layout. They take minutes and they are the difference
between a structure that survives the second board and one that has to be refactored.

\begin{enumerate}
    \item \textbf{How many chips?} If more than one, the HAL is per chip (or per family, when
          the parts are close enough that one driver set serves them). This is what
          \codeinline{nrf-hal} does per chip and \codeinline{stm32f4xx-hal} does per family;
          \codeinline{esp-hal} instead keeps one crate and selects the chip by feature, which
          works when the vendor's crates already share their register code.
    \item \textbf{How many boards, and can one sit on top of another?} Boards that are
          alternatives become one BSP crate each. Boards that \emph{stack} become one crate
          each \emph{plus} a one-way dependency: the carrier depends on the module. A single
          crate that knows both cannot answer ``does this work without the carrier?''.
    \item \textbf{Are there variants of the same board?} A variant — same PCB, different
          parts fitted — is a \emph{feature}, not a crate. A separate crate duplicates the
          board init and everything the variants share, in exchange for a name that repeats
          what the feature flag already says.
    \item \textbf{Is the wiring obvious or hard-won?} Obvious wiring (an LED on a documented
          pin) may live in an example. Hard-won wiring — a swapped bus, a power gate, a pin
          that the vendor's own map gets wrong — belongs in a BSP, once, with the measurement
          recorded next to it. This is the question that decides between Layout A and
          Layout C for most hobby-scale projects.
\end{enumerate}

The same questions decide where an algorithm goes, with one addition: \textbf{if it can be
tested on a host, it must be}, so it cannot share a crate with anything that needs a target
triple.

\subsection{Adding Support for a New Chip or Board: the Order of Work}

Having decided the layout, the order below is what this project would follow again. Each step
produces something checkable before the next one starts.

\begin{enumerate}
    \item \textbf{Registers first.} Obtain the SVD, generate the PAC, and build one example
          that only writes a register and reads it back. If the SVD is wrong, everything
          downstream inherits the error, and a five-line example finds that in minutes.
    \item \textbf{Clock and pins.} Prove that the peripheral clock gate is on and that a GPIO
          changes state. This is the step where a missing
          \codeinline{CMU\_CLKEN0} bit appears as a bus fault at the first register access.
    \item \textbf{One transport.} Bring up the simplest bus the board has, and read a device
          identifier through it. A chip identifier or a WHO\_AM\_I is a free proof that
          addressing, timing and the byte order are all correct.
    \item \textbf{Drivers, in the HAL, with pins as arguments.} Nothing in this layer may name
          a port. If a driver needs a pin, take it as a parameter; the BSP will supply it.
    \item \textbf{The BSP for one board, and its examples.} Put the wiring in one place, in
          the crate named after the board whose documentation describes the feature, and
          write the example that proves the board works end to end.
    \item \textbf{A second board, if there is one, and the dependency direction.} Add the
          carrier crate only after the module crate stands alone, and make the dependency
          point one way. If it cannot, the split is wrong.
    \item \textbf{The gate.} Before more examples are written, make the machine check what
          the review would check: zero warnings, every example builds, every domain crate
          tests on the host, and no domain crate depends on the hardware stack.
\end{enumerate}

\begin{warningbox}
    \textbf{The step people skip is the first one.} Generating the PAC and immediately
    writing drivers feels productive, and it hides an SVD error until the driver is
    complicated enough that the error looks like a driver bug. The five-line register
    read-back is the cheapest test in the whole project.
\end{warningbox}

\subsection{What This Costs}

Honesty about a design decision includes its price:

\begin{itemize}
    \item \textbf{More manifests.} Two BSP crates mean two \filepath{Cargo.toml} files, two
          example folders and a dependency between them that has to be kept one-way.
    \item \textbf{Some knowledge is duplicated on purpose.} The clock setup lives in the
          module crate and the base depends on the module rather than re-initialising; a
          project with two unrelated boards could not share it this way.
    \item \textbf{The split is a judgement, not a theorem.} ``Which plastic is this pin on?''
          is unambiguous for an LED and less so for a Grove connector that reaches the base
          through the module's castellated pads. The rule that resolves it: the BSP is named
          after the board whose \emph{documentation} describes the feature.
\end{itemize}

\begin{notebox}
    \textbf{The test to apply to your own board.} Ask which of your examples would still make
    sense to someone holding only the chip. If the answer is ``none of them'', the board has
    swallowed the crate and a split is due. If the answer is ``all of them'', you may not need
    a BSP at all — Layout A is not a mistake, it is a choice with a cost that is sometimes
    worth paying.
\end{notebox}
